\documentclass[a4paper]{article}
\usepackage{graphicx}
\usepackage{amsmath,amssymb}
\usepackage{hyperref}
\usepackage{minted}

\title{Jahr}
\date{}
\begin{document}
    \section{Nur das Jahr }
    \maketitle
    In German, years are usually pronounced as full numbers.
    
    For years from 1100 to 1999, the most common pattern is:
    \[
        \text{<century> hundred <rest>}
    \]
    Examples:
    \begin{itemize}
        \item 1998: \textit{neunzehnhundertachtundneunzig}
        \item 1984: \textit{neunzehnhundertvierundachtzig}
        \item 1975: \textit{neunzehnhundertfünfundsiebzig}
    \end{itemize}
    
    For 2000 and later, the standard form is:
    \[
        \text{zweitausend + <rest>}
    \]
    Examples:
    \begin{itemize}
        \item 2005: \textit{zweitausendfünf}
        \item 2012: \textit{zweitausendzwölf}
        \item 2024: \textit{zweitausendvierundzwanzig}
    \end{itemize}
    
    In casual speech, some people also say:
    \[
        2024:\ \textit{zwanzig vierundzwanzig}
    \]
    but the neutral/standard form remains \textit{zweitausendvierundzwanzig}.

\end{document}